
\chapter{مقدمه}
\label{chap:introduction}

این فصل مسئله‌ی مرکزی پایان‌نامه را معرفی می‌کند و چارچوبی را می‌سازد که فصل‌های بعدی بر آن استوارند. نخست توضیح داده می‌شود چرا افت دقت بخش‌بندی خودکار در فاز بدون‌کنتراست سی‌تی شکمی مسئله‌ای بالینی و نه صرفاً فنی است؛ سپس این مسئله به چهار زیرمسئله‌ی مشخص تجزیه می‌شود که ستون‌بندی مشترک فصل‌های کارهای پیشین، روش پیشنهادی و آزمایش‌ها را تعیین می‌کنند. در ادامه، چالش‌های فنی و داده‌ای که حل این زیرمسئله‌ها را دشوار می‌کنند، اهداف پژوهش، و ساختار کلی پایان‌نامه ارائه می‌شود.

\section{اهمیت مسئله}
\label{sec:importance}

سی‌تی شکمی\footnote{\lr{Abdominal CT}} یکی از پرکاربردترین وجهیت‌های تصویربرداری در تشخیص، مرحله‌بندی و پیگیری بیماری‌های کبد، کلیه، پانکراس و سایر اندام‌های شکمی است. در عمل بالینی، این تصویربرداری معمولاً در چند فاز کنتراست\footnote{\lr{Contrast phase}} انجام می‌شود: نخست یک اسکن بدون‌کنتراست\footnote{\lr{Non-contrast (NC) phase}}، پیش از تزریق ماده‌ی حاجب\footnote{\lr{Contrast media}}، و سپس یک یا چند اسکن کنتراست‌دار\footnote{\lr{Contrast-enhanced (CE) phase}} در فازهای شریانی، وریدی یا تأخیری. تشخیص بسیاری از ضایعات کانونی کبد، به‌ویژه کارسینوم سلول کبدی\footnote{\lr{Hepatocellular carcinoma (HCC)}}، مستقیماً به الگوی افزایش کنتراست و شست‌وشوی\footnote{\lr{Washout}} آن میان فازهای مختلف وابسته است~\cite{easl2018,forner2018}؛ الگویی که تنها با ترکیب دقیق اطلاعات چند فاز قابل‌مشاهده است.

با این حال، فاز بدون‌کنتراست کاربردی صرفاً تکمیلی ندارد. در بسیاری از پروتکل‌های بالینی، از جمله تروما، غربالگری فرصت‌طلبانه\footnote{\lr{Opportunistic screening}}، و بیمارانی که به دلیل نارسایی کلیوی یا سابقه‌ی حساسیت، امکان تزریق ماده‌ی حاجب را ندارند، فاز بدون‌کنتراست تنها تصویر در دسترس است. بخش‌بندی خودکار اندام‌های شکمی روی همین فاز، زیرساخت اندازه‌گیری حجمی، برنامه‌ریزی درمان و ردیابی روند بیماری در چنین بیمارانی است.

اما همان‌طور که در فصل~\ref{chap:literature} به‌تفصیل مرور می‌شود، دقت مدل‌های بخش‌بندی روی فاز بدون‌کنتراست، به دلیل کاهش نسبت کنتراست‌به‌نویز\footnote{\lr{Contrast-to-noise ratio (CNR)}} میان بافت‌های نرم مجاور، به‌طور نظام‌مند و وابسته به اندام، پایین‌تر از فاز کنتراست‌دار است~\cite{jeon2024fully,abel2025intraindividual,Zho24}. نتیجه آن است که دقیقاً همان گروه از بیماران که کمترین دسترسی را به تصویربرداری کنتراست‌دار دارند، از ضعیف‌ترین خروجی بخش‌بندی خودکار بهره می‌برند. این افت، به‌جای آن‌که محدودیتی ذاتی و اجتناب‌ناپذیر تلقی شود، نقطه‌ی شروع این پایان‌نامه است.

\section{تعریف مسئله}
\label{sec:problem-definition}

هدف کلی این پایان‌نامه، کاهش شکاف دقت بخش‌بندی چندعضوی میان فاز بدون‌کنتراست و فاز کنتراست‌دار در سی‌تی شکمی است، از طریق دو راهبرد مکمل: هم‌ترازسازی هندسی فازهای موجود از طریق ثبت تصویر، و بازسازی مستقیم اطلاعات فاز کنتراست‌دار از روی فاز بدون‌کنتراست از طریق سنتز. هر دو راهبرد بر زیرساختی مشترک تکیه دارند که پیش از ورود به آن‌ها باید مستحکم شود: شناخت دقیق محدودیت‌های مجموعه‌داده‌های موجود، و شناسایی صحیح فاز کنتراست هر بررسی.

بر این اساس، مسئله‌ی کلی به چهار زیرمسئله تجزیه می‌شود که در ادامه‌ی این بخش معرفی می‌شوند. همین چهار زیرمسئله، با همین ترتیب، ستون‌بندی فصل‌های~\ref{chap:literature} تا~\ref{chap:experiments} را نیز تعیین می‌کنند، به‌گونه‌ای که خواننده بتواند مسیر «مسئله، پیشینه، روش، آزمایش» را برای هر زیرمسئله به‌صورت مستقل دنبال کند.

\subsection{مرور انتقادی مجموعه‌داده‌های سی‌تی شکمی}
\label{subsec:sp-datasets}

پیش از طراحی هر راه‌حل، این پروژه با یک مرور نظام‌مند بر \عدد{۴۵} مجموعه‌داده‌ی عمومی سی‌تی شکمی آغاز شد که وضعیت اکوسیستم داده‌ی موجود را از منظر بازاستفاده، سوگیری جمعیت‌شناختی و جغرافیایی، و به‌طور مشخص کمیابی داده‌ی جفت‌شده\footnote{\lr{Paired}} میان فاز بدون‌کنتراست و کنتراست‌دار مستند می‌کند~\cite{danaei2026abdomenctreview}. یافته‌های این مرور در بخش~\ref{sec:datasets} به‌تفصیل ارائه می‌شود و مستقیماً انگیزه‌ی دو زیرمسئله‌ی پایانی را شکل می‌دهد: وقتی داده‌ی جفت‌شده‌ی طبیعی در دسترس نیست، یا باید فازهای موجود را پس‌پردازشی هم‌تراز کرد، یا فاز غایب را از روی فاز موجود بازسازی کرد.

\subsection{تشخیص فاز کنتراست}
\label{subsec:sp-phase}

بهره‌برداری از هر داده‌ی چندفازی، اعم از عمومی یا خام گردآوری‌شده در این پروژه، به شناسایی صحیح فاز هر سری تصویر وابسته است. از آن‌جا که برچسب‌های سرآیند \lr{DICOM} در عمل ناقص یا ناسازگارند (بخش~\ref{sec:phase-detection})، در این پروژه یک مدل تشخیص فاز کنتراست بر روی داده‌ی خام گردآوری‌شده توسعه و ارزیابی می‌شود. این مدل، که جزئیات آن در فصل~\ref{chap:method} و نتایج آن در فصل~\ref{chap:experiments} ارائه می‌شود، زیرساخت عملیاتی دو زیرمسئله‌ی بعدی است.

\subsection{ثبت تصویر بین‌فازی}
\label{subsec:sp-registration}

این پروژه یک پایپ‌لاین ثبت تصویر چندگذره\footnote{\lr{Multi-pass}} برای هم‌ترازسازی هندسی فاز بدون‌کنتراست با فاز کنتراست‌دار طراحی و ارزیابی می‌کند. معیارهای کلاسیک اطلاعات متقابل در سطح کل تصویر تعریف می‌شوند و از همین رو در برابر تغییر ناهمگن کنتراست شکننده‌اند (بخش~\ref{sec:registration}). در مقابل، پایپ‌لاین پیشنهادی از مقداردهی اولیه‌ی صلب مبتنی‌بر نشانه‌های استخوانی، گذرهای ساختاری و اطلاعات متقابل محدودشده به نقاب اندام، و پالایش نهایی غیرصلب بهره می‌گیرد تا دامنه‌ی محاسبه‌ی معیار تراز را به نواحی پایدار در برابر فاز کنتراست محدود کند. این پایپ‌لاین روی مجموعه‌داده‌ی عمومی \lr{VinDr-Multiphase}~\cite{vindr} ارزیابی شده است.

\subsection{سنتز فاز کنتراست‌دار از فاز بدون‌کنتراست}
\label{subsec:sp-synthesis}

جفت‌های هم‌ترازشده‌ی حاصل از این ثبت تصویر سپس برای آموزش و ارزیابی مدل‌های سنتز فاز کنتراست‌دار از روی فاز بدون‌کنتراست به کار می‌روند (بخش~\ref{sec:synthesis}). در نهایت، در این پروژه بررسی می‌شود که آیا ترکیب هم‌ترازسازی و سنتز بهبودیافته، شکاف دقت بخش‌بندی چندعضوی میان دو فاز را که در فصل~\ref{chap:literature} مستند می‌شود، کاهش می‌دهد.

\section{چالش‌ها}
\label{sec:challenges}

پیگیری این چهار زیرمسئله در طول پروژه با چند چالش فنی مشترک همراه بود که هر یک، انتخاب‌های روشی فصل~\ref{chap:method} را شکل داده‌اند.

نخستین چالش، ناسازگاری فراداده است. برچسب‌های فاز کنتراست در سرآیند \lr{DICOM} داده‌ی خام این پروژه، مشابه بسیاری از مجموعه‌داده‌های عمومی، ناقص یا ناسازگارند و امکان تکیه‌ی مستقیم بر آن‌ها برای آموزش یا ارزیابی وجود ندارد. همین محدودیت است که تشخیص فاز از روی محتوای تصویر را از یک گام جانبی به یک پیش‌نیاز تبدیل می‌کند.

چالش دوم به ماهیت خود تزریق ماده‌ی حاجب بازمی‌گردد: تزریق، رابطه‌ی آماری میان شدت روشنایی دو فاز را به‌شکل ناهمگن در فضا تغییر می‌دهد. عروق و پارانشیم پرخون روشن‌تر می‌شوند، در حالی که استخوان و چربی تقریباً بدون تغییر می‌مانند. این ناهمگنی، فرض بنیادین معیارهای کلاسیک ثبت تصویر مبتنی‌بر اطلاعات متقابل را نقض می‌کند و طراحی معیار تراز مقاوم را ضروری می‌سازد.

چالش سوم، کمیابی داده است. جفت طبیعی و هم‌زمان دو فاز برای یک بیمار به‌ندرت در دسترس است، و همین کمبود، هم ارزیابی دقیق ثبت تصویر و هم آموزش نظارت‌شده‌ی مدل‌های سنتز را دشوار می‌کند. سرانجام، ادبیات پژوهشی مرتبط، به‌ویژه در حوزه‌ی سنتز \lr{NC}-به-\lr{CE}، از ناهمگنی معیار ارزیابی و مجموعه‌داده رنج می‌برد؛ ناهمگنی‌ای که مقایسه‌ی مستقیم نتایج و انتخاب یک مبنای معتبر را دشوار می‌کند (بخش~\ref{sec:synthesis}).

\section{اهداف پژوهش}
\label{sec:goals}

هدف کلی این پژوهش آن است که نشان دهد شکاف دقت بخش‌بندی چندعضوی میان فاز بدون‌کنتراست و کنتراست‌دار، با مداخله در سطح داده و پیش‌پردازش، و بدون بازطراحی خود مدل بخش‌بندی، تا چه اندازه قابل کاهش است. برای رسیدن به این هدف کلی، چهار زیرمسئله‌ی بخش~\ref{sec:problem-definition} به چهار هدف عملیاتی ترجمه می‌شوند و یک هدف پنجم، اثر تجمیعی آن‌ها را می‌سنجد:

\begin{itemize}
\item مستندسازی نظام‌مند وضعیت سوگیری و کفایت مجموعه‌داده‌های عمومی سی‌تی شکمی، با تمرکز ویژه بر کمیابی داده‌ی جفت‌شده‌ی چندفازی، تا دامنه‌ی واقعی داده‌ی در دسترس برای دو هدف بعدی روشن شود.
\item توسعه و ارزیابی مدلی برای تشخیص قابل‌اعتماد فاز کنتراست از روی محتوای خام تصویر، مستقل از برچسب‌های ناقص سرآیند \lr{DICOM}، به‌عنوان زیرساخت عملیاتی بهره‌برداری از داده‌ی چندفازی.
\item طراحی یک پایپ‌لاین ثبت تصویر چندگذره که دامنه‌ی محاسبه‌ی معیار تراز را به نواحی پایدار در برابر تغییر فاز کنتراست محدود می‌کند، و ارزیابی مقایسه‌ای آن در برابر رویکردهای کلاسیک مبتنی‌بر اطلاعات متقابل.
\item بهره‌گیری از جفت‌های هم‌ترازشده‌ی حاصل برای آموزش و ارزیابی مدل‌های سنتز فاز کنتراست‌دار از روی فاز بدون‌کنتراست.
\item سنجش این‌که آیا ترکیب این مداخلات، شکاف دقت بخش‌بندی چندعضوی میان فاز بدون‌کنتراست و کنتراست‌دار را کاهش می‌دهد.
\end{itemize}

\section{ساختار پایان‌نامه}
\label{sec:thesis-structure}

باقی این پایان‌نامه به شرح زیر سازمان‌دهی شده است. فصل~\ref{chap:literature} پس از یک پیش‌زمینه‌ی کوتاه بالینی درباره‌ی سی‌تی چندفازی، پیشینه‌ی پژوهشی پنج حوزه‌ی مرتبط را مرور می‌کند: بخش‌بندی چندعضوی و شکاف کیفیت فاز بدون‌کنتراست، که مسئله‌ی چتر این پایان‌نامه را تعریف می‌کند، و پس از آن چهار حوزه‌ی متناظر با چهار زیرمسئله‌ی بخش~\ref{sec:problem-definition}؛ فصل با جمع‌بندی جای خالی مشترک میان آن‌ها به پایان می‌رسد. فصل~\ref{chap:method} روش پیشنهادی این پایان‌نامه، شامل جزئیات مدل تشخیص فاز، پایپ‌لاین ثبت تصویر چندگذره، و مدل‌های سنتز \lr{NC}-به-\lr{CE} را ارائه می‌دهد. فصل~\ref{chap:experiments} نتایج تجربی هر یک از این سه مؤلفه، و نیز ارزیابی نهایی اثر آن‌ها بر شکاف دقت بخش‌بندی چندعضوی را گزارش می‌کند. سرانجام، فصل~\ref{chap:conclusions} یافته‌های پایان‌نامه را جمع‌بندی کرده و کاستی‌ها و مسیرهای پژوهشی آتی را مطرح می‌کند.

\section{جمع‌بندی}
\label{sec:intro-summary}

فاز بدون‌کنتراست سی‌تی شکمی، با وجود کاربرد بالینی گسترده و نقش آن به‌عنوان تنها تصویر در دسترس برای گروهی از بیماران، همچنان از افت نظام‌مند دقت بخش‌بندی چندعضوی رنج می‌برد. استدلال این فصل آن بود که ریشه‌ی این افت را نباید تنها در معماری مدل‌ها جست‌وجو کرد: کاهش کنتراست بافتی، ناسازگاری فراداده‌ی فاز، و کمیابی داده‌ی جفت‌شده‌ی چندفازی، سه عامل مستقلی‌اند که همگی پیش از مرحله‌ی مدل‌سازی عمل می‌کنند. بر همین مبنا، مسئله به چهار زیرمسئله تجزیه شد که پاسخ به آن‌ها، به‌ترتیب، در فصل‌های بعدی دنبال می‌شود.
